% !TEX program = xelatex
% 海口经济学院本科毕业论文（设计）LaTeX 模板
% 依据：教务处《附件13 论文格式（修订10.21）》与《附件12 装订封面》
% 编译：xelatex -> bibtex -> xelatex -> xelatex
\documentclass[12pt,a4paper,fontset=none]{ctexart}

%====================== 页面与字体 ======================
\usepackage{geometry}
\geometry{top=30mm,bottom=30mm,left=25mm,right=20mm}
\usepackage{fontspec}
\usepackage{xeCJK}
% MIT; default preview uses TeX-distributed fonts, NOT school typography.
\RequirePackage{fontspec,xeCJK,etoolbox}
\providecommand{\ThesisFontMode}{preview}
\ifdefstring{\ThesisFontMode}{school}{%
  \newcommand{\ThesisCheckFont}[1]{\IfFontExistsTF{#1}{}{\PackageError{thesis-fonts}{Missing font: #1}{Install legally licensed fonts or use THESIS_FONTS=preview.}}}
  \ThesisCheckFont{Times New Roman}\ThesisCheckFont{SimSun}
  \ThesisCheckFont{SimHei}\ThesisCheckFont{FangSong}\ThesisCheckFont{KaiTi}
  \setmainfont{Times New Roman}
  \setCJKmainfont[AutoFakeBold=2]{SimSun}
  \setCJKsansfont[AutoFakeBold=2]{SimHei}
  \setCJKmonofont{FangSong}
  \setCJKfamilyfont{zhsong}[AutoFakeBold=2]{SimSun}
  \setCJKfamilyfont{zhhei}[AutoFakeBold=2]{SimHei}
  \setCJKfamilyfont{zhkai}[AutoFakeBold=2]{KaiTi}
  \setCJKfamilyfont{zhfs}{FangSong}
}{%
  \ifdefstring{\ThesisFontMode}{preview}{}{\PackageError{thesis-fonts}{Invalid font mode}{Use preview or school.}}
  \setmainfont{texgyretermes-regular.otf}[BoldFont=texgyretermes-bold.otf,ItalicFont=texgyretermes-italic.otf,BoldItalicFont=texgyretermes-bolditalic.otf]
  \setCJKmainfont{FandolSong-Regular.otf}[BoldFont=FandolSong-Bold.otf]
  \setCJKsansfont{FandolHei-Regular.otf}[BoldFont=FandolHei-Bold.otf]
  \setCJKmonofont{FandolFang-Regular.otf}
  \setCJKfamilyfont{zhsong}{FandolSong-Regular.otf}[BoldFont=FandolSong-Bold.otf]
  \setCJKfamilyfont{zhhei}{FandolHei-Regular.otf}[BoldFont=FandolHei-Bold.otf]
  \setCJKfamilyfont{zhkai}{FandolKai-Regular.otf}[AutoFakeBold=2]
  \setCJKfamilyfont{zhfs}{FandolFang-Regular.otf}
}
% Geometric/shapes may not exist in the serif font; ctex's open symbol font.
\xeCJKDeclareCharClass{CJK}{"2460 -> "24FF,"3200 -> "32FF}

\newcommand{\song}{\CJKfamily{zhsong}}
\newcommand{\hei}{\CJKfamily{zhhei}}
\newcommand{\kai}{\CJKfamily{zhkai}}
\newcommand{\fs}{\CJKfamily{zhfs}}
% 可换行字段，禁止使用makebox把长题目悄悄画到页面外。
\newcommand{\coverfield}[1]{\parbox[c]{9.8cm}{\centering #1\par\kern2pt\hrule height0.4pt}}

\usepackage[fleqn]{amsmath}
\usepackage{amssymb}
\usepackage{graphicx}
\usepackage{booktabs}
\usepackage{etoolbox}
\usepackage{caption}
\usepackage{titlesec}
\usepackage{tocloft}
\usepackage{fancyhdr}
\usepackage[hidelinks]{hyperref}
\usepackage[super]{gbt7714}
\usepackage{indentfirst}
\usepackage{tabularx,longtable,listings,perpage,tikz,pgfplots}
\usetikzlibrary{arrows.meta,positioning}
\pgfplotsset{compat=1.18}
\xeCJKDeclareCharClass{CJK}{"2460 -> "24FF}
\MakePerPage{footnote}
\renewcommand{\thefootnote}{\ifcase\value{footnote}\or①\or②\or③\or④\or⑤\or⑥\or⑦\or⑧\or⑨\or⑩\or⑪\or⑫\or⑬\or⑭\or⑮\or⑯\or⑰\or⑱\or⑲\or⑳\else\arabic{footnote}\fi}
\lstset{basicstyle=\ttfamily\small,breaklines=true,columns=fullflexible,frame=single,numbers=left,numberstyle=\tiny,showstringspaces=false}

%====================== 论文信息（改这里） ======================
\newcommand{\thesistitle}{基于深度学习的图像识别方法研究}
\newcommand{\thesistitleeng}{Research on Image Recognition Based on Deep Learning}
\newcommand{\thesisauthor}{张\quad 三}
\newcommand{\thesisid}{2022000001}
\newcommand{\thesisschool}{信息工程学院}
\newcommand{\thesismajor}{计算机科学与技术}
\newcommand{\thesisclass}{2022 级 1 班}
\newcommand{\thesisadvisor}{李\quad 四}
\newcommand{\thesisdate}{二〇二六年六月}

%====================== 标题样式（正文三级标题） ======================
\setcounter{secnumdepth}{3}
\setcounter{tocdepth}{3}
\renewcommand{\thesection}{\arabic{section}}
\renewcommand{\thesubsection}{\thesection.\arabic{subsection}}
\renewcommand{\thesubsubsection}{\thesubsection.\arabic{subsubsection}}

\titleformat{\section}{\hei\zihao{-3}}{\thesection}{0.6em}{}
\titleformat{\subsection}{\hei\zihao{4}}{\thesubsection}{0.6em}{}
\titleformat{\subsubsection}{\hei\zihao{-4}}{\thesubsubsection}{0.6em}{}
\titlespacing*{\section}{0pt}{0.6em}{0.4em}
\titlespacing*{\subsection}{0pt}{0.4em}{0.3em}
\titlespacing*{\subsubsection}{0pt}{0.3em}{0.2em}

%====================== 目录样式 ======================
\renewcommand{\contentsname}{目\hspace{2em}录}
\renewcommand{\cfttoctitlefont}{\hfill\song\zihao{-2}\bfseries}
\renewcommand{\cftaftertoctitle}{\hfill\mbox{}}
\renewcommand{\cftsecfont}{\song\zihao{-4}}
\renewcommand{\cftsecpagefont}{\song\zihao{-4}}
\renewcommand{\cftsubsecfont}{\song\zihao{-4}}
\renewcommand{\cftsubsecpagefont}{\song\zihao{-4}}
\renewcommand{\cftsubsubsecfont}{\song\zihao{-4}}
\renewcommand{\cftsubsubsecpagefont}{\song\zihao{-4}}
\renewcommand{\cftdotsep}{\cftnodots}
\setlength{\cftbeforesecskip}{0pt}
\setlength{\cftbeforesubsecskip}{0pt}
\setlength{\cftbeforesubsubsecskip}{0pt}
\setlength{\cftsecindent}{0pt}
\setlength{\cftsubsecindent}{1.5em}
\setlength{\cftsubsubsecindent}{3em}

%====================== 页眉页脚 ======================
\fancypagestyle{main}{%
  \fancyhf{}%
  \fancyfoot[C]{\song\zihao{5}\thepage}%
  \renewcommand{\headrulewidth}{0pt}%
}
\pagestyle{main}
\tocloftpagestyle{empty}

%====================== 图、表、公式 ======================
\DeclareCaptionFont{wu}{\song\zihao{5}}
\DeclareCaptionFont{wuhei}{\hei\zihao{5}}
\captionsetup{font=wuhei,labelsep=space,justification=centering}
\newcommand{\hkctableunit}{}
\newcommand{\tableunit}[1]{\renewcommand{\hkctableunit}{#1}}
% 表号左、标题中、计量单位右；三栏均允许换行。
\DeclareCaptionFormat{hkctable}{%
  \parbox[t]{.20\linewidth}{\raggedright #1}%
  \parbox[t]{.60\linewidth}{\centering #3}%
  \parbox[t]{.20\linewidth}{\raggedleft\hkctableunit}\par}
\AtBeginEnvironment{table}{\song\zihao{5}\renewcommand{\hkctableunit}{}}
\setlength{\heavyrulewidth}{1.5pt}
\setlength{\lightrulewidth}{0.5pt}
\captionsetup[table]{position=top,format=hkctable,singlelinecheck=false}
\renewcommand{\figurename}{图}
\renewcommand{\tablename}{表}
\renewcommand{\thefigure}{\arabic{figure}}
\renewcommand{\thetable}{\arabic{table}}
\renewcommand{\theequation}{\arabic{equation}}
\makeatletter
\renewcommand{\tagform@}[1]{\maketag@@@{（#1）式}}
\makeatother
\setlength{\mathindent}{2em}

\begin{document}
\linespread{1.5}\zihao{-4}\selectfont
\setlength{\parindent}{2em}
\setlength{\parskip}{0pt}

%====================== 封面（附件13） ======================
\begin{titlepage}
\thispagestyle{empty}
\begin{center}
\vspace*{1.0cm}
{\hei\zihao{1}\bfseries 海口经济学院}\\[1.5cm]
{\hei\zihao{2} 本\hspace{0.30em}科\hspace{0.30em}生\hspace{0.30em}毕\hspace{0.30em}业\hspace{0.30em}论\hspace{0.30em}文（设计）}\\[2.2cm]
\renewcommand{\arraystretch}{2.0}
\begin{tabular}{r@{\hspace{0.3em}}l}
{\hei\zihao{4} 题\quad 目：} & \coverfield{\hei\zihao{3}\thesistitle} \\[0.2cm]
{\hei\zihao{4} 姓\quad 名：} & \coverfield{\kai\zihao{-3}\thesisauthor} \\[0.2cm]
{\hei\zihao{4} 学\quad 号：} & \coverfield{\kai\zihao{-3}\thesisid} \\[0.2cm]
{\hei\zihao{4} 所在学院：} & \coverfield{\kai\zihao{-3}\thesisschool} \\[0.2cm]
{\hei\zihao{4} 专业班级：} & \coverfield{\kai\zihao{-3}\thesismajor\ \thesisclass} \\[0.2cm]
{\hei\zihao{4} 指导教师：} & \coverfield{\kai\zihao{-3}\thesisadvisor} \\
\end{tabular}
\end{center}

\vfill
\begin{center}
{\hei\zihao{3}\thesisdate}
\end{center}
\end{titlepage}

%====================== 声明 ======================
\newpage
\thispagestyle{empty}
\begin{center}
{\hei\zihao{2} 声\quad 明}
\end{center}
\vspace{0.8em}
{\song\zihao{-4}\hspace*{2em}
此处为原创性声明占位，不是学校声明原文。正式提交请自行取得所属学院当届声明，按要求填写并由本人确认、签名；本项目不附官方表单，也不代签。本页说明须在提交前替换，不可当作正式声明。

\vspace{3.2cm}
\hfill 本人签名：\underline{\makebox[3.2cm]{}}\par
\vspace{1.2em}
\hfill 日\quad 期：\underline{\makebox[3.2cm]{}}\ 年\ \underline{\makebox[1.2cm]{}}\ 月\ \underline{\makebox[1.2cm]{}}\ 日\par
}

%====================== 目录 ======================
\newpage
\thispagestyle{empty}
\begingroup
\titleformat{\section}[block]{\song\zihao{-2}\bfseries\centering}{}{0pt}{}
\tableofcontents
\thispagestyle{empty}
\endgroup

%====================== 摘要（页码从 1 开始） ======================
\newpage
\pagenumbering{arabic}
\pagestyle{main}
\addcontentsline{toc}{section}{摘要}
\addcontentsline{toc}{section}{关键词}

\begin{center}
{\hei\zihao{-2}\bfseries \thesistitle}
\end{center}
\vspace{0.8em}

\noindent{\kai\zihao{-4}\bfseries 摘\quad 要：}{\kai\zihao{5}
本文为图像识别研究的写作与排版示例，不报告已完成的模型训练或实验发现。模板围绕问题定义、相关工作、方法描述、实验方案、结果呈现和研究局限组织正文，提供可编辑的流程图、公式、三线表、曲线图、脚注及程序清单，并示范参考文献和附录的使用。文中出现的准确率等数值均为合成排版数据，不能据此声称任何方法优于基线。正式写作应替换为本人开展的研究，报告数据来源、划分方式、评价条件和真实结果，避免将测试集用于调参或把重复测试误当独立实验。使用者还应核对所在学院当年规定、签名要求和电子文档格式；本文件能够生成PDF，但不自动满足学校要求提交Word的条款。}

\vspace{0.5em}
\noindent{\kai\zihao{-4}\bfseries 关键词：}{\kai\zihao{5}深度学习；图像识别；卷积神经网络；注意力机制}

\vspace{1.4em}
\begin{center}
{\rmfamily\zihao{-3}\bfseries \thesistitleeng}
\end{center}
\vspace{0.6em}

\noindent{\rmfamily\zihao{5}\bfseries Abstract:}\ {\rmfamily\zihao{5}
This document is a writing and typesetting example for an image recognition thesis. It demonstrates research questions, methods, experimental design, tables, figures, equations, notes, and references. No model was trained for this document, and all numerical results are synthetic. They do not support claims of improved accuracy. Replace all examples with independently conducted research and verify current college submission requirements.}

\vspace{0.5em}
\noindent{\rmfamily\zihao{5}\bfseries Key words:}\ {\rmfamily\zihao{5}deep learning; image recognition; convolutional neural network; attention mechanism}

%====================== 正文 ======================
\newpage
% 学校特殊排序由本地hkcu-gbt7714.bst完成；不再手写nocite排序。
\section*{引言}\addcontentsline{toc}{section}{引言}
图像识别是一个具有明确输入和输出的研究任务，但研究价值不能只由选用的网络名称决定。本示例展示如何把研究问题、方法、实验和结论组织成可复核的论文，不提供真实实验成果。文中数值和图形均供排版测试，不能作为技术结论引用。
\section{研究背景与问题定义}
\subsection{研究对象与任务边界}
首先给出图像来源、类别定义、使用场景和评价目标。明确研究是在封闭数据集上比较算法，还是针对某个实际场景开展部署。样本规模、类别分布和标注质量都应说明，不能用“公开数据集”四字替代来源、版本和许可。
\subsection{相关研究与理论基础}
综述应比较已有工作的假设、优势及不足。可参考真实教材及原始论文\cite{zhou2016,vaswani2017,lecun2015}，但不能仅凭引用这些资料就声称掌握或改进了对应方法。请在实际阅读后按研究问题组织综述，并说明本文与已有工作的关系。
\subsubsection{研究目标与评价约束}
将目标转化为可观测指标，同时约束计算成本、数据使用范围和错误风险。准确率只是指标之一；类别不平衡时，还需考虑每类表现、混淆矩阵或其他合适指标。本脚注只补充解释\footnote{注释为解释性材料；参考文献用于标注来源。官方要求脚注①②…每页重置。}，不代替正式引文。
\section{方法设计与实现}
\subsection{处理流程与接口}
图~\ref{fig:pipeline} 展示一种研究流程的组织方式。训练、验证和测试应分开管理；归一化参数及特征选择只能在允许的数据范围内估计，防止信息泄漏。
\begin{figure}[htbp]\centering
\begin{tikzpicture}[node distance=5mm,box/.style={draw,minimum width=68mm,minimum height=10mm,align=center,font=\small},>=Stealth]
\node[box](a){数据来源核验与划分};
\node[box,below=of a](b){仅使用训练数据拟合模型};
\node[box,below=of b](c){验证集选择参数};
\node[box,below=of c](d){锁定方案后报告测试结果};
\draw[->](a)--(b);\draw[->](b)--(c);\draw[->](c)--(d);
\end{tikzpicture}
\caption{实验组织流程示例}\label{fig:pipeline}
\end{figure}
\subsection{数学描述与符号定义}
设样本的真实类别为 $y_i$，模型预测为 $\hat{y}_i$，准确率可写为式~\eqref{eq:accuracy}。公式编号示例采用官方“（1）式”形式，正文引用时不重复添加“式”字。
\begin{equation}\label{eq:accuracy}
\operatorname{Accuracy}=\frac{1}{n}\sum_{i=1}^{n}\mathbf{1}(y_i=\hat{y}_i).
\end{equation}
式中，$n$ 为样本数，$\mathbf{1}(\cdot)$ 为示性函数。对同一测试集反复挑选最优方案会使测试结果偏乐观，因此评价程序也应作为方法的一部分报告。
\subsection{实现细节与复现说明}
记录运行环境、依赖版本、随机种子、训练轮数、优化器、学习率策略及模型选择标准。正文解释关键改动，完整代码和数据说明宜放入附录或受控的归档位置。不同方案的比较应保持数据划分和训练预算一致。
\section{实验方案与结果呈现}
\subsection{实验设置与对照}
表~\ref{tab:protocol} 为实验设置表的编写示例。所有条件均须在真实研究中补齐；空缺信息不能靠模板推断。
\begin{table}[htbp]\centering\caption{实验设置记录项}\label{tab:protocol}
\begin{tabularx}{\linewidth}{lX}\toprule 项目 & 应报告的内容\\\midrule
数据与许可 & 名称、版本、来源、访问日期及允许用途\\
数据划分 & 训练/验证/测试划分规则，是否按个体分组\\
模型设置 & 网络配置、初始化、训练预算及参数选择依据\\
评价方式 & 指标、重复次数、置信区间及异常运行处理\\\bottomrule\end{tabularx}
\end{table}
\subsection{合成数据的展示}
表~\ref{tab:synthetic} 和图~\ref{fig:synthetic} 仅为排版示例。它们不是实验结果，不能得出改进方法有效的结论。表题在上、图题在下；单位写入表头或坐标轴，图表需在正文中先被引出。
\begin{table}[htbp]\centering\tableunit{单位见列题}\caption{合成数据示例（非模型测试结果）}\label{tab:synthetic}
\begin{tabularx}{\linewidth}{@{}>{\raggedright\arraybackslash}X *{3}{>{\raggedleft\arraybackslash}X}@{}}
\toprule 配置 & 轮次/次 & 示例指标/\% & 示例时间/s\\\midrule
A &10&80.0&12.0\\B&20&83.0&24.0\\C&30&84.0&36.0\\\bottomrule\end{tabularx}
\par\smallskip{\zihao{-5}数据来源：人为构造，仅供版式检查。}
\end{table}
\begin{figure}[htbp]\centering\begin{tikzpicture}
\begin{axis}[width=.85\linewidth,height=55mm,xlabel={轮次/次},ylabel={示例指标/\%},grid=major,xmin=8,xmax=32,ymin=78,ymax=86]
\addplot[black,mark=*]coordinates{(10,80)(20,83)(30,84)};
\end{axis}\end{tikzpicture}
\caption{合成曲线示例（非训练或测试记录）}\label{fig:synthetic}
\end{figure}
\subsection{误差分析与消融设计}
真实结果分析应检查失败样本、类别差异和数据偏差。消融实验每次只改变明确的因素，并报告多次运行的离散程度；不能只选择最好的一次结果。统计检验应满足其适用前提，不应为得到显著性反复尝试而不作说明。
\section{讨论与适用范围}
\subsection{结果解释}
将“观察到的结果”“可能的解释”和“尚未验证的推测”分开表达。改进幅度需要放在误差范围、数据规模和计算成本中解释，不能以单一数据点宣称普适性。
\subsection{局限与实际应用}
说明数据覆盖不足、标注误差和跨场景迁移的风险。涉及人脸、医疗图像或其他敏感数据时，应遵守适用的隐私、许可和伦理要求。合成示例不支持任何应用安全性判断。
\subsection{格式与归档核查}
确认各级标题、摘要、关键词、参考文献和致谢已进入目录，所有交叉引用均可解析。官方附件要求电子文档为 Word，因此 LaTeX 生成的 PDF 仅是排版成果，不等于自动取得提交资格；请提前向学院确认。

\section{结论}
本模板展示了图像识别研究论文中主要排版元素及其引用方式，不提供真实模型性能或实验结论。正式结论应回应研究问题，明确证据支持的发现、适用范围和局限，不把未来计划写成已完成成果。

%====================== 参考文献 ======================
\par\vspace{3\baselineskip}
\begingroup
\renewcommand{\bibsection}{%
  \phantomsection
  \addcontentsline{toc}{section}{参考文献}%
  \noindent{\hei\zihao{-4}参考文献}\par
  \vspace{0.8em}%
}
\renewcommand{\bibfont}{\song\zihao{5}\linespread{1}\selectfont}
\bibliographystyle{hkcu-gbt7714}
\bibliography{refs}
\endgroup

\clearpage
\section*{附录（可选）}\addcontentsline{toc}{section}{附录（可选）}
\subsection*{附录 A\quad 复现实验清单}
\begin{longtable}{lp{.68\linewidth}}\toprule 材料 & 记录说明\\\midrule\endhead
环境 & 操作系统、软件版本、依赖与运行命令\\
数据 & 来源、版本、划分文件和许可\\
模型 & 参数配置、随机种子和权重归档\\
分析 & 原始输出、绘图脚本和异常记录\\\bottomrule\end{longtable}
\subsection*{附录 B\quad 指标计算示例}
\begin{lstlisting}[language=Python]
# Synthetic labels: not an experimental result.
truth = [0, 1, 1, 0]
prediction = [0, 1, 0, 0]
accuracy = sum(a == b for a, b in zip(truth, prediction)) / len(truth)
print(accuracy)
\end{lstlisting}


%====================== 致谢 ======================
\newpage
\section*{致\quad 谢}
\addcontentsline{toc}{section}{致谢}

感谢指导教师在论文选题、研究方案和写作过程中给予的悉心指导，感谢同学和家人在学习期间的支持与帮助。

\end{document}
